\documentclass[11pt]{article}

\usepackage[a4paper,margin=1in]{geometry}
\usepackage{amsmath}
\usepackage{amssymb}
\usepackage{booktabs}
\usepackage{enumitem}
\usepackage{graphicx}
\usepackage[hyphens]{url}
\usepackage{hyperref}
\usepackage{microtype}
\usepackage{tabularx}
\usepackage{xcolor}
\usepackage[numbers]{natbib}

\hypersetup{
    colorlinks=true,
    linkcolor=blue,
    citecolor=blue,
    urlcolor=blue
}

\title{Toward Skill Evolution for Blind Vulnerability Analysis Agents}
\author{Authors omitted in working draft}
\date{August 2026}

\begin{document}
\maketitle

\begin{abstract}
Large language model agents are increasingly equipped with reusable natural-language
skills, but their role in vulnerability analysis remains under-specified.  In this
domain, agents repeatedly perform similar procedural steps: locate an entry point,
identify a dispatch branch, trace attacker-controlled data, reach a sink, and decide
whether the path represents a real vulnerability.  These recurring steps suggest that
blind vulnerability analysis should benefit from reusable skills.  At the same time,
our early experiments show that static skill libraries are insufficient: a skill that
helps one firmware family may be irrelevant or harmful on another, and apparently
successful runs may still be poor candidates for long-term reuse.

This paper argues that \emph{skill evolution for blind vulnerability analysis} is an
important research problem in its own right.  The key question is not merely how to
write or retrieve skills, but how to decide when the result of a completed blind run
should be converted into a reusable skill that future analyses will depend on.  We
therefore present a first framework-oriented formulation of this lifecycle: blind remote
analysis, server-side skill selection and injection, local external confirmation,
run-level feedback construction, candidate skill generation, and gated publication into
a live library.

Our current contribution is intentionally scoped as a first arXiv-stage paper.  We do
not claim that stable autonomous skill improvement has already been achieved, nor that
the current system can discover unknown vulnerabilities end-to-end.  Instead, we aim to
establish the problem, the lifecycle, and the initial evidence that make this line of
work worth pursuing.  Empirically, leakage-controlled reruns and frozen multi-case
firmware studies show that skill usage can materially change vulnerability localization
outcomes, but the effect is heterogeneous, family-dependent, and not yet sufficient to
prove strict necessity across all cases.  Archived gate outcomes further suggest that
candidate generation is not the only difficulty: deciding whether a candidate skill
should enter the live library is itself a central open problem.  We present the current
prototype and preliminary evidence as a foundation for future work on attribution,
publication control, and discovery-oriented evaluation.
\end{abstract}

\noindent\textbf{Keywords:} LLM agents; skill evolution; blind vulnerability analysis;
firmware security; external confirmation; reusable procedural knowledge.

\section{Introduction}
\label{sec:intro}

Large language model (LLM) agents increasingly rely on reusable memories, natural-language
instructions, and explicit skill libraries to improve over repeated tasks.  Systems such
as Voyager~\cite{voyager}, ExpeL~\cite{expel}, and SkillClaw~\cite{skillclaw} show that
agents can accumulate procedural knowledge and reuse it later.  Other recent work studies
context-to-skill induction~\cite{ctx2skill}, on-policy skill distillation~\cite{opid},
environment-grounded skill refinement~\cite{spark,skillevolbench}, and verifier-guided
skill revision~\cite{trace2skill,veriskill}.  Together, these directions suggest that
skill is becoming a practical interface between one-shot model behavior and longer-lived
agentic adaptation.

Vulnerability analysis is a particularly natural setting for this idea.  In blind
analysis or discovery-oriented auditing, the agent repeatedly faces variations of the
same procedural problem: identify where attacker-controlled input enters, decide which
branch or handler deserves attention, trace data or control to a dangerous operation,
and determine whether the behavior is exploitable.  These patterns recur across firmware
families, binaries, and vulnerability types.  If any domain should benefit from reusable
skills, vulnerability analysis is one of them.

However, the presence of repeated structure does not imply that a fixed skill library is
enough.  Recent work has shown that skills can also be harmful: a topically relevant
skill can push an agent toward the wrong abstraction level or the wrong execution
path~\cite{skillharmful}.  Our own experiments show the same pattern in firmware
analysis.  A skill that improves vulnerability localization on one family may hurt
another; a more specific skill can outperform a generic one in some cases and underperform
it in others.  This means the real question is not simply whether vulnerability-analysis
agents need skills, but why those skills must evolve from real runs instead of remaining
static templates.

That requirement creates a second, harder problem.  A completed run may appear
successful, yet still be a bad source for long-term skill reuse.  It may be correct only
by chance, correct but redundant because the model would have succeeded anyway, or
apparently correct while drifting to the wrong vulnerable function within the same
firmware family.  Thus, once we move from ``should agents use skills?'' to ``should
agents evolve skills?'', a new research object appears:

\begin{quote}
When should the result of a completed blind vulnerability-analysis run be converted into
a reusable skill that will influence future analyses?
\end{quote}

This paper focuses on the combination of \emph{agent-based vulnerability analysis} and
\emph{skill evolution}, rather than on either topic in isolation.  We do not claim that
the full problem is already solved.  Instead, we argue that the path itself remains
under-explored: blind vulnerability analysis benefits from skills; static skills are not
enough; and converting run outcomes into evolving skills introduces new questions of
evidence, attribution, and publication control.

Our prototype setting couples remote blind analysis with local external confirmation.
An agent operates in an isolated analysis environment; the local system archives the
run, scores the final answer, checks it against external evidence, and only then allows
an evolution service to propose candidate skill updates.  The resulting workflow is not
yet a mature autonomous discovery system.  It is, however, a concrete and testable
prototype for studying whether the outputs of blind vulnerability-analysis runs should
become persistent skills.

This paper makes four contributions:
\begin{enumerate}[leftmargin=*]
    \item \textbf{We identify skill evolution for blind vulnerability analysis as a
    distinct research problem.}  The core question is not merely how to write or retrieve
    skills, but how to move from a completed run to a reusable live skill without
    polluting future analyses.

    \item \textbf{We propose a prototype lifecycle for that problem.}  The lifecycle
    connects blind remote analysis, server-side skill injection, local external
    confirmation, run-level feedback, candidate skill generation, and gated publication.

    \item \textbf{We present initial empirical evidence that skill matters but is
    heterogeneous.}  Clean reruns and frozen firmware studies show that skill can
    materially change vulnerability localization, while family-specific experiments show
    that more specific skills are not uniformly better.

    \item \textbf{We argue that candidate evaluation and publication are central future
    challenges.}  Current gate results suggest that generating candidate skills is not
    enough; the harder question is deciding which candidates deserve to enter the live
    library.
\end{enumerate}

The scope of the paper is therefore deliberate: this is a framework paper with
preliminary evidence, written to establish the research direction early rather than to
claim that stable autonomous skill evolution has already been achieved.

\section{Why Skill, Why Evolution, Why Here}
\label{sec:why}

\subsection{Why vulnerability analysis needs skills}

Blind vulnerability analysis is not a single prediction problem.  It is a repeated
procedural task.  In firmware and binary auditing, analysts frequently perform the same
types of reasoning steps:
\begin{itemize}[leftmargin=*]
    \item find the entry point or handler that exposes attacker-controlled input;
    \item identify the relevant branch, mode, or dispatch path;
    \item follow dataflow toward dangerous operations or implicit sinks;
    \item distinguish the true target path from nearby but wrong high-risk paths;
    \item convert local evidence into a vulnerability-level conclusion.
\end{itemize}

These steps are too structured to be treated as fully one-off reasoning.  A reusable
skill can encode how to traverse such a path more reliably than a bare prompt that
starts from scratch every time.

\subsection{Why static skills are not enough}

A fixed library is appealing, but our experiments suggest it is not sufficient.  In some
firmware families, a narrower family-specific skill improves localization by steering the
model toward the right handler.  In others, the same narrowing hurts because the family
contains dense and overlapping handler clusters.  Likewise, a skill that is topically
relevant can still be procedurally wrong: it may talk about command injection in
general, yet encourage the model to summarize a category of sinks rather than narrow to
the correct branch and target function.

This means a vulnerability-analysis skill library cannot be treated as static reference
text.  It must be revised in response to real runs, real failures, and real case
families.

\subsection{Why skill evolution is harder than skill retrieval}

Once we allow skill revision, we face a stronger requirement than ordinary task success.
Suppose a run gets the right answer.  That does not necessarily mean its selected skills
were useful, or that the resulting procedural trace should be promoted into future runs.
At least four cases are possible:
\begin{enumerate}[leftmargin=*]
    \item the skill was genuinely useful and should be retained;
    \item the skill was harmless but unnecessary;
    \item the run succeeded for reasons unrelated to the skill;
    \item the run found a dangerous but wrong target path and should not be reinforced.
\end{enumerate}

Skill evolution therefore raises a governance problem, not just a generation problem.

\section{Problem Statement}
\label{sec:problem}

We model a run as
\begin{equation}
    r = (x, S_r, \tau, y, E),
\end{equation}
where $x$ is a blind vulnerability-analysis task, $S_r$ is the set of selected skills,
$\tau$ is the trajectory, $y$ is the final answer, and $E$ is the externally collected
evidence.

The paper is centered on three questions:
\begin{enumerate}[leftmargin=*]
    \item \textbf{Skill effect:} how much do selected skills change the final analysis
    outcome?
    \item \textbf{Run-to-skill transition:} when does a completed run contain reusable
    procedural knowledge?
    \item \textbf{Publication control:} when should a candidate update be allowed to
    influence later runs through the live skill library?
\end{enumerate}

The second and third questions are the ones we believe are still under-articulated in
the current literature on security agents and skill libraries.

To make this concrete, let a completed run $r$ produce a candidate skill update
\begin{equation}
    \Delta s = \mathcal{C}(r),
\end{equation}
where $\mathcal{C}$ is a candidate-construction procedure over the archived run.  A
separate publication policy then decides whether $\Delta s$ should influence later runs:
\begin{equation}
    \mathcal{P}(r, \Delta s, E) \in \{\texttt{publish}, \texttt{hold}, \texttt{reject}\}.
\end{equation}
In principle, $\mathcal{P}$ should depend not only on final answer quality, but also on
whether the run reached the right target path, whether the resulting procedural pattern
is reusable, and whether promoting it would risk polluting later analyses.  Our current
prototype only approximates this policy, but the equation makes the intended research
object explicit: skill evolution is not only about generating $\Delta s$, but also about
studying the decision rule $\mathcal{P}$.

For the current prototype, $\mathcal{P}$ is operationalized through a replay-style gate.
Given a baseline run $r$ and a candidate skill update $\Delta s$, the system constructs
a follow-up run
\begin{equation}
    \hat{r} = \mathcal{R}(x, S_r \cup \{\Delta s\}),
\end{equation}
where $\mathcal{R}$ re-executes the same task under the candidate-augmented skill set.
The candidate is then evaluated under a conservative publication rule of the form
\begin{equation}
    \mathcal{P}(r, \Delta s, E) = \texttt{publish}
    \iff
    \bigl[m(\hat{r}) \geq m(r) - \varepsilon \bigr]
    \wedge
    \bigl[c(\hat{r},E)=\texttt{pass}\bigr],
\end{equation}
where $m(\cdot)$ is a task-level quality signal and $c(\cdot)$ checks externally
confirmed consistency.  This replay gate is intentionally conservative: it is a
practical approximation to publication control, not a complete solution to causal skill
attribution.

\section{Prototype Framework}
\label{sec:framework}

\subsection{Split-plane execution}

The framework separates two planes:
\begin{itemize}[leftmargin=*]
    \item an \textbf{analysis plane}, where the agent performs blind vulnerability
    analysis in a remote environment; and
    \item a \textbf{confirmation plane}, where the local system scores outputs, gathers
    external evidence, archives sessions, and drives skill evolution.
\end{itemize}

This separation reduces direct answer leakage while still allowing the system to process
evidence after the run.

\subsection{Skill injection and session archival}

Selected skills are injected server-side, and each run is archived together with its
selected skill set and immutable session segment.  This does not solve causal
attribution, but it makes later analysis auditable: we can at least determine what
procedural prior the model saw.

\subsection{External confirmation}

The current prototype supports multiple evidence forms:
\begin{itemize}[leftmargin=*]
    \item source-backed predicates;
    \item logic harnesses;
    \item behavioral oracles; and
    \item sanitizer-backed runtime evidence.
\end{itemize}

These signals are not treated as proofs of skill contribution.  Instead, they constrain
what kinds of run outcomes are eligible to enter the evolution loop.

\subsection{Feedback, candidates, and gate}

After scoring and confirmation, the system constructs run-level feedback, generates
candidate skill updates, and stages them behind a gate before publication.  This gate is
the current boundary between \emph{interesting candidate text} and \emph{trusted live
skill}.

\section{Prototype Realization}
\label{sec:prototype}

The current implementation extends SkillClaw with:
\begin{itemize}[leftmargin=*]
    \item remote blind execution support;
    \item firmware- and source-oriented benchmark case descriptions;
    \item server-side skill selection with inline and catalog-style injection paths;
    \item target-aware scoring and diagnostic flags;
    \item run archival and handoff records;
    \item candidate skill generation through an evolution service; and
    \item validation-gated publication from candidate storage into a live skill
    directory.
\end{itemize}

The implementation is best viewed as a research prototype rather than a finalized product.
Its value lies in making the run-to-skill transition measurable and debuggable.

Operationally, the system maintains a distinction between a \emph{live} skill library
used for subsequent analysis, a \emph{candidate} area that stores evolution outputs
awaiting review, and a \emph{shared/archive} area that records historical skill states
and session-linked artifacts.  This separation is not only an engineering convenience.
It is necessary if publication is itself treated as a research problem rather than as an
automatic side effect of candidate generation.

\section{Preliminary Experimental Evidence}
\label{sec:evidence}

\subsection{What the current experiments are meant to show}

The present experiments are not designed to prove that autonomous skill evolution is
already reliable.  They are designed to answer a narrower set of questions:
\begin{enumerate}[leftmargin=*]
    \item does skill materially change blind vulnerability-analysis outcomes?
    \item are those effects uniform, or highly case-dependent?
    \item does the evolution loop face its main difficulty in candidate generation or in
    candidate publication?
\end{enumerate}

The reporting policy in this draft follows the same principle.  Clean reruns, the
frozen firmware protocol, family-specific distinguishing studies, and archived gate
records are treated as \emph{primary evidence}.  Earlier development-time runs remain
useful for diagnosis, but are not used as headline quantitative results once answer
leakage, overly permissive scoring, or protocol contamination were identified.

\begin{table}[t]
\centering
\caption{Current experimental suites used in this draft.}
\label{tab:arxiv_experiment_suites}
\scriptsize
\begin{tabular}{p{0.16\textwidth} p{0.18\textwidth} p{0.18\textwidth} p{0.38\textwidth}}
\toprule
Study & Target & Conditions & Primary role in this paper \\
\midrule
Leakage-controlled clean rerun & 84 source/firmware reruns & no-skill / weak-skill / oracle-skill & Re-check whether earlier skill gains survive after removing answer-bearing prompt content \\
Frozen firmware necessity study & 153 runs over 17 firmware targets & no-skill / weak-skill / oracle-skill & Measure whether stronger skills change true-target localization and decoy attraction under a fixed protocol \\
Family-specific distinguishing studies & Belkin F9K1122; Tenda FH451 & generic oracle / distinguishing oracle & Test whether narrower family-specific skills help only when the family exposes usable discriminative cues \\
Gate archive analysis & 94 archived gate decisions & candidate rerun + gate decision logs & Examine whether the main bottleneck already lies in candidate publication rather than candidate generation \\
\bottomrule
\end{tabular}
\end{table}

\subsection{Leakage-controlled reruns}

An 84-run three-condition rerun was used to clean earlier overly optimistic oracle-style
results.  After removing answer-bearing skill content, oracle-skill still outperformed
weak-skill and no-skill, while sharply reducing decoy hits.  This shows that skill
effects remain real after leakage cleanup, rather than collapsing entirely.

More specifically, under this cleaned protocol, oracle-skill achieved a $35.7\%$ true
hit rate, compared with $14.3\%$ for weak-skill and $3.6\%$ for no-skill.  Just as
importantly, decoy hits dropped from 9 and 12 down to 1.  This matters because a large
part of firmware-analysis failure is not ``finding nothing,'' but being pulled toward a
wrong yet still dangerous handler that appears superficially plausible.

\subsection{Frozen 153-run firmware study}

A larger frozen protocol was then run across 17 firmware cases under three conditions:
no-skill, weak-skill, and oracle-skill.  Under the primary reporting protocol, oracle
skill achieved $46.7\%$ true hits, versus $11.1\%$ for weak-skill and $6.7\%$ for
no-skill, while reducing decoy hits from 28/45 and 24/45 to 2/45.  These numbers are
not yet sufficient to prove strict necessity for every case, but they do show that skill
materially changes localization quality.

The practical implication is that stronger skill formulation changes the agent's target
selection behavior, not merely its narrative quality.  Several families that were
previously dominated by nearby decoy handlers move toward the intended vulnerable path
once the agent is given a more structured procedural prior.  At the same time, a large
residual ``NO'' region remains, which is why the current paper avoids claiming solved
necessity or solved discovery.

\subsection{Family-specific distinguishing studies}

The current results also show that specificity is not uniformly beneficial.  In the
Belkin F9K1122 family, a distinguishing skill increased accuracy from $20.0\%$ to
$66.7\%$.  In the Tenda FH451 family, the same idea reduced accuracy from $40.0\%$ to
$26.7\%$.  This suggests that the usefulness of a refined skill depends on whether the
family is structurally separable rather than merely topically similar.

This contrast is important for the paper's main argument.  If narrower skills always
helped, then skill evolution would mostly be a content-accumulation problem.  Instead,
the current evidence suggests that the \emph{shape} of a skill must match the structure
of the target family.  In some families, a distinguishing cue is enough to dislodge a
persistent attractor.  In others, the same style of refinement only overcommits the
model to unstable cues.

\subsection{Gate behavior}

Archived gate records show that candidate generation is not the only challenge.  The
system can already produce candidates and stage them for follow-up validation, but it
still lacks a strong criterion for deciding which candidates genuinely deserve to enter
the live library.  Replay-only gates can reject obviously poor candidates, yet they do
not by themselves prove that a published skill contributed something necessary rather
than merely harmless.

Across 94 archived gate decisions, 12 were published and 82 were rejected.  The
important point is not the raw acceptance rate.  The important point is that publication
already behaves like a distinct filtering problem with its own failure modes, operating
after candidate text has been generated.  In other words, the main unsolved issue in the
current prototype is no longer ``can the system write candidate skill text?'' but
``under what evidence should that text be trusted enough to enter the live library?''

The deeper limitation is that a replay-only gate cannot cleanly separate three cases:
genuine improvement, harmless redundancy, and wrong-path success that still survives a
permissive metric.  This is why the present paper treats gate design as a research
problem rather than as a solved engineering component.

\begin{table}[t]
\centering
\caption{Condensed summary of the current main empirical signals.}
\label{tab:arxiv_main_results}
\footnotesize
\begin{tabular}{p{0.28\textwidth} p{0.17\textwidth} p{0.17\textwidth} p{0.28\textwidth}}
\toprule
Study & Weaker condition & Stronger condition & What it shows \\
\midrule
Clean rerun (84 runs) & no-skill: 3.6\%; weak-skill: 14.3\% & oracle-skill: 35.7\% & Skill effect remains after leakage cleanup; decoys collapse from 9/12 to 1 \\
Frozen firmware study (153 runs) & no-skill: 6.7\%; weak-skill: 11.1\% & oracle-skill: 46.7\% & Stronger skills change target localization, not only output wording \\
Belkin F9K1122 & generic oracle: 20.0\% & distinguishing oracle: 66.7\% & Family-specific refinement can recover the right target when separable cues exist \\
Tenda FH451 & generic oracle: 40.0\% & distinguishing oracle: 26.7\% & Narrower skills are not uniformly better across dense handler families \\
Gate archive (94 decisions) & 82 rejected & 12 published & Candidate generation is possible; publication remains the harder filtering problem \\
\bottomrule
\end{tabular}
\end{table}

\section{Current Evidence and Scope Boundaries}
\label{sec:status}

\subsection{What the current work already supports}

The present prototype and experiments already support the following claims:
\begin{itemize}[leftmargin=*]
    \item blind vulnerability analysis is a plausible setting for skill-based agents;
    \item skill usage can materially alter vulnerability localization results;
    \item the effect of skill is heterogeneous and family-dependent;
    \item a full run-to-candidate-to-gate loop can be prototyped and studied.
\end{itemize}

\subsection{What remains unresolved}

The current work does \emph{not} yet prove:
\begin{itemize}[leftmargin=*]
    \item stable autonomous skill improvement over long horizons;
    \item strict necessity of a particular skill for all target cases;
    \item robust counterfactual attribution of skill contribution; or
    \item end-to-end unknown-vulnerability discovery in a fully autonomous sense.
\end{itemize}

This boundary is intentional.  The first goal of this draft is to establish that the
research path exists and is worth occupying now, not to over-claim that all of its hard
subproblems are already solved.

\subsection{What this draft is claiming}

The intended claim of this paper is modest but concrete.  We claim that
\emph{skill evolution for blind vulnerability analysis} is a distinct and
under-explored research direction; that a complete run-to-skill lifecycle can already be
prototyped and studied in a realistic firmware setting; and that early evidence is
sufficient to justify treating candidate evaluation and publication as first-class
research problems.  We do not claim that the present system has solved skill evolution
or autonomous vulnerability discovery.

\section{Research Agenda}
\label{sec:agenda}

We view three follow-up directions as central.

\paragraph{Candidate evaluation.}
The next step is to determine whether a candidate skill merely preserves performance or
actually improves later runs under controlled comparisons.

\paragraph{Publication policy.}
Future versions of the framework should move beyond replay-only gates toward stronger
counterfactual or retention-aware publication policies.

\paragraph{Discovery-oriented evaluation.}
The benchmark should broaden from current firmware and confirmed-vulnerability settings
toward more open-ended discovery tasks, while preserving leakage control and external
confirmation.

\section{Related Work}
\label{sec:related}

\subsection{Skill libraries, induction, and accumulation}

Voyager~\cite{voyager}, ExpeL~\cite{expel}, Reflexion-style trajectory reuse, and
SkillClaw~\cite{skillclaw} together establish the broader context in which LLM agents
accumulate reusable procedural knowledge over time.  Ctx2Skill~\cite{ctx2skill} and
OPID~\cite{opid} show that skills can be induced from context or distilled from richer
policies.  These works motivate the present paper's starting assumption: skill is a
reasonable substrate for making agent behavior persistent across tasks.

\subsection{Verifier- and environment-grounded skill refinement}

Recent work increasingly grounds skill refinement in environment evidence or external
verification rather than in unconstrained self-reflection alone.  SPARK~\cite{spark}
emphasizes evidence-backed refinement, Trace2Skill~\cite{trace2skill} studies
trajectory-to-skill conversion with verifier support, VeriSkill~\cite{veriskill}
specializes this idea to program-verification tasks, and SkillEvolBench~\cite{skillevolbench}
turns episodic-to-procedural evolution into an explicit evaluation problem.  Our work is
aligned with this trend, but shifts the focus to a security-specific question: how a
blind vulnerability-analysis run should be archived, confirmed, and filtered before it
can shape a live skill library.

\subsection{Skill failure, attribution, and governance}

Skill-induced failure is now a recognized problem.  Skill Harmful~\cite{skillharmful}
shows that apparently relevant skills can degrade performance by steering the model
toward the wrong procedure.  Misevolution~\cite{misevolution} highlights that repeated
self-improvement can preserve poor or unsafe behaviors as persistent skills.  This line
of work is directly relevant to our setting, because vulnerability analysis is unusually
sensitive to wrong-but-plausible paths: a candidate skill can look useful while
systematically biasing the model toward a nearby dangerous decoy.  The present paper
therefore treats publication control as a first-class research problem rather than as a
simple engineering afterthought.

\subsection{Security agents, playbooks, and realistic benchmarks}

On the security side, PentestGPT~\cite{pentestgpt}, empirical studies of LLM-based
vulnerability analysis~\cite{steenhoek2024vulndetect}, EvoHunt~\cite{evohunt}, and SEC-bench Pro~\cite{secbenchpro}
address different parts of the security-agent landscape: interactive offensive
workflows, vulnerability reasoning, playbook evolution, and realistic long-horizon
evaluation.  Our paper sits closest to the intersection of these threads.  Like
EvoHunt, we care about reusable procedural knowledge in security workflows.  Like
SEC-bench Pro, we care about evaluation realism and leakage control.  The difference is
that we center the run-to-library transition itself: blind execution, external
confirmation, candidate staging, and live-skill admission are the primary objects of
study.

\section{Threats to Validity}
\label{sec:threats}

\paragraph{Sample scale.}
The current experiments are larger and cleaner than earlier development-time studies, but
they are still modest in scale compared with mature security benchmarks.  Our claims are
therefore directional rather than universal.

\paragraph{Single primary model family.}
The main runs were conducted under a fixed upstream model configuration.  Skill effects
may look different under models with different priors, tool-use habits, or instruction
following behavior.

\paragraph{Static-analysis-heavy evaluation.}
Much of the current firmware evidence comes from blind static reasoning rather than from
fully executable live exploitation.  This is sufficient for studying target selection,
but weaker than a broader dynamic validation setting.

\paragraph{Attribution precision.}
The present gate and feedback signals can constrain what should not become a live skill,
but they do not yet provide a full causal decomposition of why a run succeeded.

\paragraph{Benchmark scope.}
The current paper still draws primarily from confirmed-vulnerability and controlled
family settings.  Open-ended unknown-vulnerability discovery remains a future step.

\section{Conclusion}
\label{sec:conclusion}

This paper argues that skill evolution for blind vulnerability analysis deserves to be
studied as a problem in its own right.  Vulnerability-analysis agents appear to need
skills because their work repeatedly follows structured procedures; static skills are not
enough because the value of a skill depends on target family, task structure, and
failure mode; and evolving skills introduces new questions about evidence, attribution,
and publication.

Our current prototype does not solve the full problem.  It does, however, make the path
concrete enough to study.  We provide a first lifecycle that connects blind execution,
external confirmation, run archival, candidate generation, and publication control; and
we provide preliminary evidence that skill usage materially changes vulnerability
localization while leaving candidate evaluation as the main unresolved bottleneck.

We therefore present this draft as an early research claim: the combination of
agent-based vulnerability analysis and skill evolution is both under-explored and
promising, and it merits a dedicated line of follow-up work.

\bibliographystyle{plainnat}
\bibliography{references_arxiv_checked}

\end{document}
